% Fragmento público generado desde propietarios íntegros.
% Residencia: 52.8.3.
% Fuente: ../incoming/radion_monografia_20260827/manuscrito/sections/09_sector_90120.tex:115-175
\subsubsection{Fonctionnelle hexade--octade de Barbero--Immirzi}

\paragraph{Séries orientées et saut d'incidence}

On définit
\begin{equation}
S_{90}(q)=\frac{q^{90}}{1-q^{270}},
\qquad
S_{120}(q)=\frac{q^{120}}{1-q^{360}},
\label{eq:barbero-series}
\end{equation}
et
\[
\Delta S_m=S_m(q_-)-S_m(q_+).
\]
La fonctionnelle d'inclusion marquée est
\begin{equation}
K_{6\to8}=12\Delta S_{90}-\Delta S_{120}.
\label{eq:barbero-núcleo}
\end{equation}

\paragraph{Vérification diophantienne du poids dodécaphasique \(12:-1\)}

La chaîne fondamentale dodécaphasique et son lecteur de retour, construits dans
\eqref{eq:l9s102:barbero-cadena-dodecafasica}--%
\eqref{eq:l9s102:barbero-peso-generado}, produisent douze termes de type
\(S_{90}\) et un terme de frontière de type \(-S_{120}\). Le couple orienté
est déterminé avant le calcul du coefficient de pôle.

Considérons \(aS_{90}-bS_{120}\). La vérification du coefficient de
\((1-q)^{-1}\) prend la forme
\begin{equation}
\frac{a}{270}-\frac{b}{360}=\frac1{24}.
\label{eq:pole-condition}
\end{equation}
Multiplier par \(1080\) donne
\[
4a-3b=45.
\]
L'unique couture du tour fondamental a une multiplicité absolue de un.
Sa vérification diophantienne correspond au contre-terme entier minimal
\(|b|=1\). Pour \(b=1\), \(a=12\) ; pour \(b=-1\), \(a=21/2\) n'est pas entier.
Par conséquent,
\begin{equation}
\boxed{a:b=12:1,\qquad aS_{90}-bS_{120}=12S_{90}-S_{120}.}
\label{eq:12minus1}
\end{equation}

\begin{proposition}[Unicité de la vérification diophantienne]
Le couple généré \((12,1)\) est l'unique solution entière de
\eqref{eq:pole-condition} avec canal principal positif et contre-terme
minimal \(|b|=1\).
\end{proposition}

La réalisation hexade--octade conserve les canaux \(90/120\)
préalablement construits ; la chaîne dodécaphasique apporte les multiplicités
et l'orientation de frontière ; le lecteur de retour attribue les séries aux
générateurs. Ces opérations conjointes produisent \(12S_{90}-S_{120}\).
L'équation \eqref{eq:pole-condition}, la positivité et la minimalité
conservent leur fonction de vérification arithmétique ultérieure.

La chaîne de Barbero est ultérieure :
\[
(A,C^*)\to(x,y)\to q_\pm
\to K_{6\to8}\to\gamma_{\mathrm{BI}},
\]
avec l'évaluation
\[
\gamma_{\mathrm{BI}}=0.274007672694459\ldots.
\]
Barbero--Immirzi ne génère pas rétroactivement \(C^*\), \(\Delta_4\) ni
\(\epsR\).

